\documentclass{article}
\usepackage[T1]{fontenc}
\usepackage[utf8]{inputenc}
\usepackage{kotex}
\usepackage{graphicx}
\usepackage[sort]{natbib}
\usepackage[margin=1in]{geometry}
\usepackage{setspace}
\doublespacing

\title{The Korea Discount: A Review of Corporate Governance Deficits and the Case for Structural Reform}
\author{Author Name}
\date{\today}

\begin{document}

\maketitle

\begin{abstract}
South Korean equities have traded at a persistent structural discount to developed and emerging market peers for at least three decades. This paper synthesizes the academic literature on the ``Korea Discount,'' drawing on 26 studies to assess the competing explanations---concentrated Chaebol ownership and tunneling, weak minority shareholder protections, and stagnant corporate growth prospects---and to evaluate the reform mechanisms proposed to resolve it. We find that while growth fundamentals and governance deficits are jointly implicated, the causal evidence from natural experiments consistently identifies the self-dealing channel as a direct and quantifiable source of value destruction. Institutional stewardship through the National Pension Service and the Stewardship Code has produced limited and conditional improvements, insufficient on their own to close the valuation gap. We support this conclusion with an original empirical contribution: a 20-year monthly panel of index-level price-to-book (P/B) ratios comparing Korea to Japan as a governance reform benchmark. Our results show that Korea's relative valuation penalty deepened precisely during the period when Japan escalated from voluntary codes to mandatory, exchange-level structural enforcement in 2023, while Korea remained on a voluntary program. The evidence points to a consistent conclusion across the literature and our own analysis: mitigating the Korea Discount requires hard legislative reform---specifically, amendments to the Commercial Act mandating broader fiduciary duties---rather than continued reliance on soft-law guidelines.
\end{abstract}

\newpage

\section{Literature Review}

\subsection{The Persistent P/B Disparity: Defining the Korea Discount}
The ``Korea Discount'' refers to the chronic undervaluation of South Korean equities relative to global peers, measured through depressed price-to-book (P/B) and price-to-earnings (P/E) ratios. Across 48 countries from 1990 to 2023, Korean firms exhibit persistently lower P/B and P/E multiples---a finding that survives controls for firm size, profitability, growth, GDP per capita, financial market development, and geopolitical risk across all sub-periods \citep{kim_korean_market_longterm}. At the firm level, the high prevalence of Korean firms with book-to-market ratios above one is more consistent with systematic mispricing than fundamental risk, and a governance factor---rather than geopolitical or purely fundamental factors---is the cross-sectional predictor most strongly associated with this pattern \citep{jung_korea_discount_book_to_market}.

The literature is divided, however, on which root causes are most likely. Direct tests of governance, shareholder payout, and growth potential hypotheses find that corporate governance has no statistically significant relationship with PBR, while variables capturing stagnant growth prospects---low R\&D expenditure, high tangible-asset intensity, and firm maturity---exhibit positive correlations with low PBR \citep{yang_korea_discount_growth_vs_governance}. Higher shareholder payout tendencies are also associated with \emph{lower} PBR, counter to the shareholder-return argument. Consistent with the growth view, for large manufacturing firms (which account for most of the KOSPI) dividend expansion has weak or even negative effects on PBR, while R\&D investment is the more significant driver \citep{lee_korea_discount_pbr}.

These findings directly challenge the governance-reform interpretation, and this is not fully resolved by the existing literature. The distinction between cross-sectional correlations and causal identification through natural experiments, reviewed below, partly explains the divergence: natural experiments consistently implicate governance as a direct value-destroyer even where cross-sectional regressions cannot detect it. The Korea Discount is likely overdetermined: weak growth fundamentals and governance failures reinforce one another, and reforms addressing only one dimension will underperform.

North Korean geopolitical risk is a commonly cited but empirically weak alternative. Cross-country studies that control for geopolitical risk explicitly find the Korea Discount survives intact, consistent with markets treating North Korean threats as a chronic, well-priced background condition. Section~4.3 formalizes this test; we reject the geopolitical explanation and focus on the governance and fundamental factors reviewed below.

\subsection{Root Causes: Concentrated Ownership, Chaebols, and Tunneling}
The structural foundation of the Korea Discount lies in Korea's Chaebol system---family-controlled conglomerates that dominate the Korean economy. Controlling families use well-established ``central firms'' to acquire low-profitability entities through pyramidal structures; central firms trade at a discount to other group firms, with shareholders anticipating value-destroying acquisitions, though some underperformance may reflect selection rather than tunneling since low-profitability firms are placed in pyramids by design \citep{almeida_chaebol_formation}. The cash flow rights disparity between voting and economic ownership is substantially larger than previously reported, as studies using only publicly listed firms miss the amplifying role of non-public affiliates \citep{kim_chaebol_cf_rights}.

When a Chaebol-affiliated firm makes an acquisition, its own stock price falls on average while the controlling shareholder benefits through value transfer to other group firms \citep{bae_tunneling_value_added}. Business group firms also make systematically larger charitable donations, coordinated to tunnel resources from public affiliates (where family cash flow rights are lower) to private ones (where family stakes are higher) \citep{bae_business_groups_tunneling}.

The most causally credible evidence exploits Korea's 1999 reform imposing stricter board structure requirements on large firms (assets above 2 trillion won) relative to mid-sized firms: large firms whose controllers had incentives to tunnel earned strong positive abnormal returns around the reform, and panel regressions over 1998--2004 confirm that better governance moderates the negative effect of related-party transactions on firm value and raises the sensitivity of firm profitability to industry profitability \citep{black_governance_firm_value}. This natural experiment isolates the self-dealing channel and provides the identification that cross-sectional governance regressions cannot.

Concentrated ownership also distorts the broader governance environment. Among Chaebol-affiliated firms, only shareholder rights protection retains explanatory power for dividend payouts; board effectiveness, audit competency, and disclosure transparency lose independent influence once ownership structure is accounted for \citep{njoku_corporate_governance_dividends}. This pattern shows that standard improvements in board composition and disclosure are insufficient when concentrated ownership remains entrenched. The structural consequences include a persistent holding company discount \citep{suh_holding_company_discount}---unique to Korea and absent from Japan and the US---and the use of cross-shareholding arrangements as management entrenchment devices \citep{park_cross_shareholdings_takeover}.

\subsection{The Limits of Institutional Stewardship in Korea}
A first wave of policy responses sought to leverage the National Pension Service (NPS) as an activist institutional shareholder. NPS shareholding improves ESG performance and accounting and market outcomes among Korean manufacturing firms, with ESG improvements serving as the transmission channel \citep{kim_nps_esg_manufacturing}. NPS disclosures of ESG engagement intentions are associated with significant firm value gains in the following year, with effects largest for firms with the weakest internal governance---indicating that institutional activism substitutes for, rather than complements, internal governance mechanisms \citep{lee_esg_nps_firm_value}.

More direct forms of activism show limited reach. Markets do not react to NPS ``Vote No'' press announcements in the short run---a result inconsistent with developed-market shareholder activism findings---though firms that subsequently improve internal governance after a vote-no campaign do show higher valuation over time \citep{kim_nps_governance_firm_value}. The short-run market non-reaction underscores that NPS activism lacks enforcement credibility absent structural reform. Stewardship Code adoption by institutional blockholders is positively correlated with earnings quality, but multivariate regressions fail to establish statistical significance, attributed to a time lag before shareholder activism translates into measurable governance changes \citep{park_stewardship_code_earnings}.

A broader challenge emerges from a cross-country study covering 58 markets from 2001 to 2018: Korea's formal shareholder rights scores are not distinctively lower than peers, and markets with the largest shareholder return increases actually saw declines in average firm value \citep{kim_korea_discount_causes}. What distinguishes Korea is the low value relevance of fundamentals---profitability and growth are incorporated into prices to a far lesser degree than in comparable markets (51st of 58 countries)---suggesting a short-term investment culture prevents long-term growth expectations from being capitalized.

Together, institutional activism within the current legal framework is insufficient to close the P/B gap absent structural reform to shareholder enforcement.

\subsection{Top-Down Policy: Japan's Reforms as the Benchmark}
Contemporary literature has converged on top-down institutional reform as the necessary complement to institutional activism, drawing heavily on Japan's staged governance transformation. Korea's 2024 ``Corporate Value-Up'' program's voluntary structure produced low firm participation and a weak market response---the KOSPI closed weaker in 2024, the opposite of Japan's experience following mandatory TSE P/B directives \citep{lee_value_up_effectiveness}. Japan's reforms proceeded through a staged progression: the 2014 Stewardship Code, the 2015 Corporate Governance Code, and ultimately the 2023 TSE exchange-level mandate requiring firms with P/B below 1.0 to disclose capital cost improvement plans---each stage ratcheting enforcement rather than relying on voluntary compliance.

The 2025 Korean Commercial Act amendment---which expanded directors' fiduciary duties from the corporation to shareholders broadly---provides direct evidence that legal reform shifts valuation expectations: firms with the highest book-to-market ratios earned significantly positive abnormal returns at both the preliminary expectation date and the parliamentary passage date, with the differential reaction larger around formal passage than around the preceding presidential election, and no significant reaction among high controlling-shareholder-ownership firms \citep{kang_commercial_act_market_reactions}. These patterns suggest markets expect the amendment to have selective rather than uniform effects---larger where governance reform has the most room to reduce expropriation. Foreign investors act as asymmetric disciplinary monitors, correcting both over-investment in over-growing firms and underperformance in under-growing ones, implying that reforms improving market accessibility would amplify market discipline through the foreign investor channel \citep{joo_foreign_investors_value_up}.

The emerging consensus is that resolving the Korea Discount requires a transition from voluntary guidelines to comprehensive structural mandates. Korea's current trajectory mirrors Japan's pre-2023 soft-law phase. Our empirical analysis corroborates this view. Korea's relative P/B penalty against Japan deepened precisely during the period in which Japan escalated to exchange-level enforcement while Korea remained on a voluntary program. Critically, the same pattern holds when Japan is replaced by MSCI EM Asia as the comparator---ruling out Japan-specific post-reform appreciation as the mechanical driver---consistent with the view that hard structural mandates, not soft codes, are the critical threshold for sustained valuation recovery.

\section{Data and Methodology}

\subsection{Data and Sample Construction}

Our empirical analysis rests on a monthly panel of index-level Price-to-Book (P/B) ratios spanning January 2004 to April 2026. The core universe comprises four benchmark indices: the Korea Composite Stock Price Index (KOSPI), the Tokyo Stock Price Index (TOPIX), the S\&P 500, and the MSCI Emerging Markets index. These four series constitute the primary study universe and support both descriptive and event-study analysis through December 2024, which we treat as the principal study endpoint consistent with the paper's pre-registered scope. A follow-on analysis window extending through April 2026 supports the Korea Value-Up event-study and robustness specifications, which require additional post-event observation months given the recency of Korea's 2024 reform milestones.

The regional robustness analysis draws on a broader set of MSCI country and regional indices---MSCI Hong Kong, MSCI Taiwan, MSCI EM Asia, MSCI China, MSCI India, and MSCI Indonesia---which serve two distinct purposes. In the Korea alternative-comparator robustness specification, MSCI EM Asia replaces TOPIX as the denominator of the Korea valuation spread, allowing us to assess whether Japan-specific post-reform appreciation mechanically contaminates our headline Korea--Japan spread estimate. In the synthetic control analysis, the same regional indices constitute the donor pool against which a counterfactual Japan is constructed.

All P/B series are measured at the aggregate index level and are available at monthly frequency. Index-level P/B ratios have several methodological advantages over the firm-level multiples that dominate the microeconometric governance literature reviewed above. Because they reflect the value-weighted composition of each market's listed equity, they integrate information across thousands of firms and smooth idiosyncratic cross-sectional noise. This property is particularly important for identifying structural valuation wedges---such as the Korea Discount---that persist across sectors, size groups, and firm-level governance variation. The tradeoff is that index-level analysis precludes controlling for firm-level observables such as profitability, leverage, or R\&D intensity; our identification strategy therefore relies on the plausible assumption that such compositional differences between KOSPI and TOPIX are approximately stable within event windows of the length we employ (12 months pre-event to 24 months post-event).

\subsection{Descriptive Analysis}

Figure 1 presents the 20-year P/B time series for KOSPI, TOPIX, S\&P 500, and MSCI EM over the 2004--2024 study window, with Japan's three primary governance reform dates annotated by vertical reference lines. This figure establishes the key stylized facts motivating our event-study design. The KOSPI--TOPIX gap, visible at the start of the sample and persistent throughout, narrows modestly during the global financial crisis and its aftermath, then resumes a downward trajectory by the mid-2010s. The most striking feature of the descriptive record is the divergent trajectory of the two indices following Japan's March 2023 TSE P/B reform mandate, a divergence that the event-study analysis formalizes.

\subsection{Primary Spread and Event-Study Design}

The primary outcome variable throughout our analysis is the KOSPI--TOPIX P/B spread: the monthly difference between the KOSPI index-level P/B ratio and the TOPIX index-level P/B ratio. Choosing Japan as the comparator is motivated by structural, geographic, and institutional similarities that make it the most credible counterfactual available. Japan and Korea are both capital-account-open, export-oriented economies with concentrated corporate ownership structures, and both listed equity markets are classified as developed by major index providers. Crucially, Japan has passed through a staged sequence of governance reforms over the past decade, culminating in a hard exchange-level mandate in 2023---precisely the reform trajectory that Korea's policymakers and the literature argue Korea must replicate. These features justify Japan as the benchmark against which Korean governance policy can be evaluated.

We use a stacked event-study design to estimate the cumulative abnormal spread around each reform event. The stacking approach, which constructs a separate cohort dataset for each event date and then pools them into a single regression, is well-suited to settings with multiple non-overlapping or partially overlapping reform events, because it allows each cohort's pre-event trend to serve as its own counterfactual baseline without constraining effects to be homogeneous across events \citep{cengiz_stacked_did_2019}.

For each event date $t_k$ ($k = 1, \ldots, K$), we construct a cohort panel by extracting all months within a symmetric window around that event. The stacking window for the Japan analysis runs from 36 months before to 24 months after each event date, providing a pre-event observation span sufficient to estimate a reliable baseline trend. Months within each cohort are indexed by their calendar-month distance from $t_k$, denoted $\tau \in [-36, +24]$, where negative values are pre-event and positive values are post-event. For the Korea analysis, the stacking window is $[-36, +20]$ months, with the post window truncated to 20 months to respect the available post-reform observation period through April 2026.

\paragraph{Abnormal Spread Construction.} Within each cohort, we estimate a pre-event linear trend in the KOSPI--TOPIX spread over the pre-period $\tau \in [-36, -1]$, fitting a simple OLS regression of the spread on $\tau$. This linear trend projection provides the expected spread for all months in $[-36, +24]$, including the post-event months. The abnormal spread for each month is then defined as the realized spread minus the expected spread. Normalizing the baseline period at $\tau = -1$ (the month immediately before the reform event) ensures that estimated effects capture deviations from the pre-event trajectory rather than level differences in the pre-period mean, and removes transient level shifts that may have accumulated at event announcement.

\paragraph{Stacked Regression and Cumulative Abnormal Returns.} We pool the cohort-level abnormal spread panels into a stacked dataset and estimate a saturated cohort-by-relative-time OLS model:
\begin{equation}
    \text{AbnSpread}_{k,\tau} = \sum_{k=1}^{K} \sum_{\tau \neq -1} \beta_{k,\tau} \cdot \mathbf{1}[\text{cohort} = k, \text{event\_rel\_time} = \tau] + \varepsilon_{k,\tau}
\end{equation}
where the base category is $\tau = -1$ for each cohort, so $\beta_{k,-1} \equiv 0$ by construction. The saturated design is fully flexible: it imposes no parametric restriction on the pattern of abnormal spreads within or across cohorts, and makes no restrictions on treatment effect heterogeneity across events. The cumulative abnormal return (CAR) for cohort $k$ through month $\tau^*$ is:
\begin{equation}
    \text{CAR}_{k,\tau^*} = \sum_{\tau=-12}^{\tau^*} \hat{\beta}_{k,\tau}
\end{equation}
where the event window for the reported CARs begins at $\tau = -12$ rather than the full stacking window start at $\tau = -36$, ensuring that CARs reflect the policy-relevant window around each reform date rather than the broader pre-event period used for baseline estimation.

Because the saturated cohort-by-relative-time design exhausts all degrees of freedom in the time dimension within each cohort, the design matrix is not full rank when observations from multiple cohorts overlap in calendar time---as they do in the Korea analysis, where three Value-Up reform dates cluster within a seven-month window in 2024. In this setting, robust sandwich standard errors are undefined for the saturated design, and we therefore report estimated coefficients and CARs as descriptive statistics rather than as estimates with formal inferential claims. This is the appropriate treatment in an exploratory context where overlapping event windows are a design feature to be disclosed rather than a source of bias to be eliminated.

\subsection{Japan Event-Study: Reform Benchmark Identification}

The Japan event-study uses three governance reform events: (i) the Financial Services Agency Stewardship Code (February 2014), which established behavioral expectations for institutional investors; (ii) the Tokyo Stock Exchange Corporate Governance Code (June 2015), which imposed formal board composition and disclosure standards on listed companies; and (iii) the TSE P/B Reform Request (March 2023), under which the TSE required all listed firms with P/B below 1.0 to disclose capital efficiency improvement plans or explain non-compliance. These three events span Japan's full governance reform arc from soft institutional guidance to hard exchange-level mandate, and they constitute the governance sequence that Korea's policy trajectory is being benchmarked against.

The direction of the spread effect depends on whether the reform in question caused TOPIX to appreciate or KOSPI to depreciate (or both) in relative terms. A persistent negative CAR in the KOSPI--TOPIX spread following a Japanese reform would indicate that Japan's valuation improved relative to Korea's---consistent with the governance channel if the reform was effective in raising TOPIX values but did not produce a symmetric response in KOSPI. A null or positive CAR would indicate no widening of the Korea Discount in response to Japan's reform milestones.

\subsection{Geopolitical Risk Falsification}

The literature review established that North Korean geopolitical risk is a commonly cited but empirically weak driver of the Korea Discount. We test this formally using the Geopolitical Risk index for Korea (GPRC\_KOR) constructed by \citet{caldara_gprc_2022}, which provides a monthly country-level measure of media-reported geopolitical threats. To operationalize the escalation hypothesis in its strongest form, we construct a binary escalation dummy equal to one when GPRC\_KOR exceeds its 75th percentile over the 2004--2024 study window---a threshold that identifies the most severe geopolitical risk episodes rather than the chronic background level.

We estimate the following cross-sectional time-series model:
\begin{equation}
    \text{KOSPI\_PB}_t = \alpha + \beta_1 \cdot \text{GPR\_Escalation}_t + \beta_2 \cdot \text{TOPIX\_PB}_t + \sum_{y} \gamma_y \cdot \mathbf{1}[\text{year} = y] + \varepsilon_t
\end{equation}
where the TOPIX P/B control absorbs common global equity market movements, and year fixed effects absorb secular trends in both series over the two-decade panel. Standard errors are HC3 heteroskedasticity-robust. The key parameter is $\beta_1$: if geopolitical risk is a meaningful driver of the Korea Discount, KOSPI P/B should be systematically lower in high-escalation months even after controlling for TOPIX co-movement and year-level trends. A precisely estimated zero would undermine the geopolitical risk hypothesis and support the governance-centric interpretation pursued throughout this paper.

We note that $\hat{\beta}_1$ is a partial identification result---the GPR escalation dummy may absorb global risk-off episodes that simultaneously depress multiple equity markets, in which case the coefficient is an upward-biased estimate of the pure North Korea threat effect. Our preferred interpretation of a near-zero coefficient is therefore that even this potentially attenuated estimate of the geopolitical risk channel is negligible.

\subsection{Synthetic Control Validation}

As an alternative identification strategy for the 2023 TSE P/B reform, we apply the synthetic control method of \citet{abadie_synthetic_control_2010} to construct a counterfactual Japan P/B trajectory. The treatment unit is the TOPIX (Japan) index, and the donor pool consists of seven MSCI regional and country indices that did not undergo an equivalent P/B governance mandate in 2023: MSCI Hong Kong, MSCI Taiwan, MSCI Emerging Markets, MSCI EM Asia, MSCI China, MSCI India, and MSCI Indonesia. South Korea (KOSPI) is excluded from the donor pool to preserve the independence of the primary KOSPI--TOPIX spread analysis. All donor countries are regional or emerging market peers with no exchange-level P/B mandate reform during the study period, satisfying the SUTVA assumption that donor units are unaffected by the Japanese treatment.

Japan's absolute P/B level lies structurally below several donor markets, creating a convex hull problem under which standard synthetic control matching on levels would concentrate all weight on the single lowest-valued donor. To resolve this, we demean each series to its pre-treatment sample mean before optimization, allowing the synthetic control weights to match on trend shape rather than levels. The pre-treatment period for weight estimation spans January 2004 to February 2023 (19 years), providing a long baseline over which global governance and ESG trends common to all markets are averaged out.

The synthetic Japan gap series---realized TOPIX P/B minus weighted synthetic TOPIX P/B---is our counterfactual outcome for the post-2023 period. We assess the reliability of the synthetic control estimate through two placebo tests. The in-time placebo assigns a fake treatment date of January 2019 and re-runs the full estimation on pre-actual-reform data, verifying that the pre-treatment fit holds and that no spurious post-``treatment'' gap emerges before the true reform. The in-space placebo iterates over each donor market, treating each as the treated unit with the remaining donors plus Japan constituting the placebo pool, and reports the distribution of placebo gap series. A Japan gap that is visually distinguishable from the placebo distribution in the post-2023 period constitutes informal evidence that the observed gap is attributable to the TSE reform rather than to structural differences between Japan and its regional peers.

\subsection{Korea Value-Up Event Study}

We apply the same stacked event-study design to three Korean governance reform milestones from Korea's 2024 ``Corporate Value-Up'' program: (i) the program launch (February 26, 2024), which announced the policy framework; (ii) the guidelines unveiling (May 2, 2024), which released implementation guidelines for voluntary firm participation; and (iii) the implementation push (August 12, 2024), which reiterated regulatory pressure for firms to file value enhancement plans. These dates are locked in \texttt{config.py} and imported by all analysis scripts to prevent any post-hoc selection of event dates.

The Korea event-study uses the same stacking window parameters as the Japan analysis ($[-36, +20]$ months), with the post window capped at 20 months given data availability through April 2026. The primary spread is KOSPI minus TOPIX P/B, preserving strict comparability with the Japan event-study. An alternative spread using KOSPI minus MSCI EM Asia P/B is reported as a robustness check: because MSCI EM Asia includes Korea as a constituent, this benchmark captures Korea's valuation relative to the broader Asian emerging market complex without depending on the Japan-specific post-2023 reform dynamics that drive the primary spread.

The three Korea reform dates cluster within a seven-month window (February, May, and August 2024), creating calendar-time overlap across cohort event windows. Under overlap, months in early 2024 simultaneously contribute to the post-event segment of the February cohort and the pre-event segment of the May cohort. We retain and explicitly annotate these overlapping observations rather than dropping them or restricting the event set, for two reasons. First, the Korea Value-Up policy milestones are analytically distinct, and collapsing them into a single event would discard informative variation in market responses across program stages. Second, since we report descriptive CARs rather than inferential estimates, the overlap does not distort the coefficient interpretation in the same way it would distort a difference-in-differences variance estimator.

\subsection{Counterfactual Projection}

Figure 4 extends the synthetic control analysis to a forward-looking scenario. Using the average monthly lift in the TOPIX--synthetic TOPIX gap during months 1--18 post-reform as an estimate of the monthly valuation increment attributable to Japan's governance mandate, we apply this lift to KOSPI's last observed P/B level (December 2024) to project Korea's hypothetical P/B trajectory over a 60-month horizon under the assumption that Korea implements an analogous exchange-level reform. Uncertainty is quantified by the pre-treatment root mean square prediction error (RMSPE) from the synthetic control estimation, which reflects the quality of fit in the pre-treatment period and provides a conservative model uncertainty band.

This projection is explicitly illustrative rather than a causal forecast. The lift estimate is drawn from Japan's experience under a specific institutional context, and Korea's governance trajectory, Chaebol composition, and investor base differ in ways that would affect the magnitude and timing of any valuation response. The projection's purpose is to calibrate the order of magnitude of potential gains under structural reform, not to provide a precise point estimate.

\subsection{Robustness Specifications}

We report four robustness checks for the Korea event-study. First, the \textit{narrow 2024 rollout} specification is the primary baseline, using the three February--August 2024 dates with a 20-month post window. Second, a \textit{window sensitivity} specification truncates the post window to 12 months, testing whether the headline result is sensitive to how much of the follow-on period is included given the limited post-event history. Third, the \textit{spaced follow-through} specification replaces the clustered 2024 dates with a temporally spread event sequence---the February 2024 program launch, the July 2025 mandatory governance disclosure expansion, and the February 2026 Value-Up dividend disclosure rule---allowing each event's effects to propagate without window overlap, though at the cost of a maximum post window of only 2 months for the latest event given the current study endpoint. Fourth, the Korea--Asia alternative comparator specification, discussed above, replaces TOPIX with MSCI EM Asia as the spread denominator.

Together, the primary result and these robustness specifications map out the sensitivity of our conclusions to comparator choice, event date selection, and post-event window length---the three most consequential design choices in an index-level event study of this kind.

\section{Results}

\subsection{Descriptive Evidence}

Table~\ref{tab:summary_stats} presents full-panel and sub-period summary statistics for the four benchmark indices over January 2004 to December 2024. Over the full sample, KOSPI trades at a mean P/B ratio of 1.18, compared with 1.35 for TOPIX, 1.78 for MSCI Emerging Markets, and 3.08 for the S\&P 500. The KOSPI--TOPIX P/B spread averages $-0.177$ P/B points over the 252 monthly observations, with a Newey-West standard error of $0.055$ and a $t$-statistic of $-3.23$. Against the broader emerging-market benchmark, the discount is substantially larger: the KOSPI--MSCI EM spread averages $-0.601$ P/B points, with a Newey-West standard error of $0.058$ and a $t$-statistic of $-10.31$. Both gaps are persistent across the two-decade panel, with 95 percent confidence intervals of $[-0.284, -0.069]$ for the TOPIX spread and $[-0.716, -0.486]$ for the MSCI EM spread, confirming that neither estimate reflects a short-lived episode or a single crisis event.

The sub-period decomposition reveals a structural deterioration in Korea's relative valuation that coincides with Japan's governance reform arc. In the pre-reform period (2004--2013), KOSPI averaged 1.36 and TOPIX averaged 1.41---a gap of only 0.05 P/B points that amounts to near-parity. During the reform era (2014--2022), spanning Japan's Stewardship Code through the years preceding the TSE P/B mandate, KOSPI fell to a mean of 1.02 while TOPIX averaged 1.29---a gap of 0.27 P/B points, more than five times its pre-reform level. In the post-TSE sub-period (2023--2024), the deterioration accelerated further: KOSPI averaged 0.93 while TOPIX averaged 1.36, a mean gap of 0.43 P/B points. Korea thus entered the post-TSE period trading below book value on average, a threshold the literature associates with chronic governance-related undervaluation, while Japan was re-rating upward. The MSCI EM comparison reinforces this picture: against an emerging-market benchmark that averaged 1.97 in the pre-reform period and 1.66 in the post-TSE period, Korea's discount has grown from approximately 0.61 to 0.73 P/B points, indicating that Korea's deterioration relative to emerging-market peers preceded and survived the Japan-specific reform dynamics. Figure~1 presents the full 20-year P/B time series, with Japan's three primary reform dates annotated.

\begin{figure}[htbp]
    \centering
    \includegraphics[width=\textwidth]{figure1_pb_comparison.png}
    \caption{Monthly P/B ratios for KOSPI, TOPIX, MSCI EM, and S\&P 500, January 2004--December 2024. Vertical dashed lines mark Japan's three governance reform milestones: the Stewardship Code (February 2014), the TSE Corporate Governance Code (June 2015), and the TSE P/B Reform Request (March 2023). The divergence between KOSPI and TOPIX following the 2023 mandate is the primary focus of the Japan event study in Section~4.2.}
    \label{fig:pb_comparison}
\end{figure}

\subsection{Japan Reform Benchmark}

Figure~2 presents the cumulative abnormal spreads from the Japan stacked event study across both spread definitions. The direction and magnitude of the CARs depend sharply on the nature of each reform event, illustrating the importance of distinguishing soft institutional guidance from hard exchange-level mandates.

\begin{figure}[htbp]
    \centering
    \begin{subfigure}[b]{0.49\textwidth}
        \centering
        \includegraphics[width=\textwidth]{figure_asia_event_study.png}
        \caption{MSCI EM Asia--TOPIX spread (primary)}
        \label{fig:japan_es_asia}
    \end{subfigure}
    \hfill
    \begin{subfigure}[b]{0.49\textwidth}
        \centering
        \includegraphics[width=\textwidth]{figure2_event_study.png}
        \caption{KOSPI--TOPIX spread (secondary)}
        \label{fig:japan_es_topix}
    \end{subfigure}
    \caption{Cumulative abnormal spreads for Japan's three governance reform events, estimated using the stacked cohort event-study design. Each line plots the CAR from $\tau = -12$ to the post-event window end for one cohort; the base period is $\tau = -1$. Panel~(a) uses the primary MSCI EM Asia--TOPIX P/B spread; Panel~(b) uses the secondary KOSPI--TOPIX P/B spread. A negative CAR indicates TOPIX appreciated relative to the comparator. The March 2023 TSE P/B Reform Request produces a persistently negative CAR in both panels, reaching $-2.58$ P/B points (EM Asia) and $-6.48$ P/B points (KOSPI) at the window end. The 2014 and 2015 soft-code events produce positive CARs, indicating no commensurate EM Asia or Korea underperformance relative to trend.}
    \label{fig:japan_es}
\end{figure}

\paragraph{Stewardship Code (February 2014).} In the primary MSCI EM Asia--TOPIX specification, the abnormal spread accumulated to $+1.595$ P/B points by $\tau = 20$: EM Asia held its relative position against Japan, and Japan's 2014 voluntary code produced no TOPIX re-rating that the regional basket failed to match. The secondary KOSPI--TOPIX spread corroborates this: CARs of $+1.322$ at $\tau = 12$ and $+4.291$ at $\tau = 24$ confirm that Korea specifically also did not underperform. Japan's first soft-law reform did not widen the Japan--Asia structural gap.

\paragraph{TSE Corporate Governance Code (June 2015).} The primary MSCI EM Asia--TOPIX CAR reached $+1.596$ at $\tau = 20$, again indicating no meaningful TOPIX re-rating relative to regional peers. The secondary KOSPI--TOPIX CARs ($+2.965$ at $\tau = 12$; $+7.829$ at $\tau = 24$) are quantitatively similar in direction. Neither early soft-law event constituted an enforceable governance shock capable of differentially re-rating TOPIX above regional peers.

\paragraph{TSE P/B Reform Request (March 2023).} The results for Japan's 2023 hard mandate are sharply different in both sign and trajectory. In the primary MSCI EM Asia--TOPIX specification, the abnormal spread declines immediately on announcement and accumulates monotonically to $-2.584$ P/B points by $\tau = 20$, with negative month-on-month contributions in all but one post-event month. The result is unambiguous: Japan's mandatory P/B reform lifted TOPIX above the entire EM Asia regional basket, not just relative to Korea. The secondary KOSPI--TOPIX specification shows an even larger divergence---$-6.481$ P/B points by $\tau = 24$, the largest in absolute value of any event in the sample---indicating that Korea's individual shortfall exceeded the regional average. The consistency of sign and monotone trajectory across both spread definitions identifies the 2023 TSE mandate as the structural inflection point separating Japan's governance reform arc from regional stagnation.

\subsection{Geopolitical Risk}

Table~\ref{tab:geo_risk} presents the results of the geopolitical risk falsification test. The coefficient on the GPR escalation dummy---defined as months in which the Caldara-Iacoviello GPRC\_KOR index exceeded its 75th percentile (0.370) over the 2004--2024 study window---is $-0.016$ P/B points, with an HC3 standard error of $0.019$, a $t$-statistic of $-0.835$, and a $p$-value of $0.405$. The null hypothesis that geopolitical escalation has no conditional association with KOSPI P/B cannot be rejected at any conventional significance level. The estimate is also negligible in economic magnitude: even taken at face value, the point estimate implies that the most severe 25 percent of geopolitical episodes are associated with a 0.016 P/B point decline in KOSPI, equivalent to less than one-tenth of the mean KOSPI--TOPIX spread.

The near-zero result is conservative: GPRC\_KOR may absorb global risk-off episodes that simultaneously depress multiple markets, so the test is biased toward finding a geopolitical effect---and finds none. The TOPIX P/B control loads at $0.629$ (SE $= 0.117$, $t = 5.40$, $p < 0.001$), confirming strong KOSPI co-movement with global equity conditions; after controlling for this common factor and year fixed effects, the escalation dummy adds nothing across 252 monthly observations. Figure~3 presents the conditional scatter. Spanning multiple nuclear tests, missile launches, and diplomatic episodes, the formal test finds no residual valuation effect attributable to geopolitical escalation, and the Korea Discount cannot be attributed to a geopolitical risk premium.

\begin{figure}[htbp]
    \centering
    \includegraphics[width=0.75\textwidth]{figure3_geo_risk.png}
    \caption{KOSPI P/B ratio by GPR escalation status, conditional on TOPIX P/B and year fixed effects. High-escalation months (GPR escalation dummy $= 1$) are defined as months in which GPRC\_KOR exceeds its 75th percentile (0.370) over 2004--2024. The OLS coefficient on the escalation dummy is $-0.016$ P/B points ($t = -0.835$, $p = 0.405$), indistinguishable from zero. The regression controls for TOPIX P/B and includes year fixed effects; $N = 252$ monthly observations. HC3 standard errors.}
    \label{fig:geo_risk}
\end{figure}

\subsection{Korea Value-Up Event Study}

Figure~4 presents the primary Korea Value-Up event-study CARs for the KOSPI--MSCI EM Asia spread; Figure~5 presents the secondary results for the KOSPI--TOPIX spread. Because the three Korea reform milestones cluster within a seven-month window (February--August 2024), the TOPIX comparison is subject to a specific confound: in 2024, TOPIX was still elevated from its 2023 re-rating before beginning to plateau, so the pre-event trend extrapolations for the May and August cohorts absorb part of Japan's plateau as an apparent Korea improvement. The MSCI EM Asia benchmark carries no equivalent mandate and is free of this artifact; it constitutes the primary evidence on Korea's governance trajectory.

\begin{figure}[htbp]
    \centering
    \includegraphics[width=\textwidth]{figure_korea_asia_event_study.png}
    \caption{Cumulative abnormal KOSPI--MSCI EM Asia P/B spreads for Korea's three Value-Up reform milestones: program launch (February 26, 2024), guidelines unveiled (May 2, 2024), and implementation push (August 12, 2024). Each cohort is normalized to zero at $\tau = -1$; CARs accumulate from $\tau = -12$. A negative CAR indicates Korea underperformed regional emerging-market peers relative to pre-event trend. All three cohorts produce persistently negative CARs: the launch reaches $-2.178$ at $\tau = 12$ and $-4.141$ at $\tau = 20$; guidelines reach $-1.804$ at $\tau = 12$; implementation push reaches $-2.121$ at $\tau = 12$. Because MSCI EM Asia carries no TSE-equivalent P/B mandate, these results are unaffected by Japan-specific TOPIX appreciation. See Figure~5 for the secondary KOSPI--TOPIX results.}
    \label{fig:korea_es_asia}
\end{figure}

\paragraph{KOSPI--MSCI EM Asia spread (primary).} Against a regional emerging-market benchmark unaffected by TSE-specific reform dynamics, all three Korea Value-Up milestones produce persistently negative CARs throughout the post-event window. For the program launch (February 26, 2024), the abnormal KOSPI--MSCI EM Asia spread falls to $-0.315$ P/B points by $\tau = 0$, to $-0.792$ by $\tau = 6$, to $-2.178$ by $\tau = 12$, and to $-4.141$ by $\tau = 20$, with negative month-on-month contributions in 16 of the 20 post-event months. The launch event is the cleanest of the three milestones because it is least contaminated by overlap with later cohorts' pre-event estimation windows. For the Guidelines Unveiled cohort (May 2, 2024), the CAR reaches $-0.179$ at $\tau = 3$, deepens to $-1.804$ at $\tau = 12$, and stands at $-2.681$ at $\tau = 20$. For the Implementation Push cohort (August 12, 2024), the CAR is $-0.193$ at $\tau = 0$, $-1.147$ at $\tau = 6$, and $-2.121$ at $\tau = 12$. Across all three milestones, Korea underperformed regional peers by between 1.8 and 4.1 P/B points at the 12-month horizon. The convergence of negative CARs across three distinct Value-Up milestones, and the absence of any sustained positive CAR, establishes that Korea's voluntary reform program did not produce a durable improvement relative to regional peers.

\begin{figure}[htbp]
    \centering
    \includegraphics[width=\textwidth]{figure_korea_event_study.pdf}
    \caption{Cumulative abnormal KOSPI--TOPIX P/B spreads for Korea's three Value-Up reform milestones (secondary analysis). The launch event (February 2024) produces a negative 12-month CAR of $-0.289$ P/B points. The guidelines and implementation push cohorts show positive CARs over the extended window ($+1.030$ and $+1.426$ at $\tau = 12$; $+3.305$ and $+5.470$ at $\tau = 20$), reflecting in part the partial plateauing of TOPIX in late 2024 after its 2023 re-rating rather than genuine Korea improvement. Compare with Figure~4 (MSCI EM Asia spread) for the unconfounded primary result.}
    \label{fig:korea_es_topix}
\end{figure}

\paragraph{KOSPI--TOPIX spread (secondary).} The TOPIX-based results illustrate the confound discussed above. For the program launch event (February 26, 2024), the abnormal KOSPI--TOPIX spread falls immediately on announcement: the CAR reaches $-0.190$ at $\tau = 0$ and $-0.289$ at $\tau = 12$, consistent with the primary specification's negative result. The Guidelines Unveiled (May 2, 2024) and Implementation Push (August 12, 2024) cohorts, however, show positive CARs accumulating to $+1.030$ and $+1.426$ respectively at $\tau = 12$, and $+3.305$ and $+5.470$ at $\tau = 20$. These apparent improvements are artifacts: the pre-event trend extrapolations for these later cohorts begin during TOPIX's 2023--2024 peak, so the subsequent TOPIX plateau is absorbed as a positive spread contribution. The primary MSCI EM Asia specification---which shows uniformly negative CARs for all three milestones---rules out any interpretation of the positive TOPIX CARs as genuine Korean valuation improvement.

Across all six event-study configurations (three milestones, two spreads), only the TOPIX specification for the two later milestones produces positive CARs, and those results are explained by the TOPIX confound identified above. The primary MSCI EM Asia evidence uniformly shows Korea losing ground to regional peers across every Value-Up milestone.

\subsection{Robustness}

We assess robustness along two dimensions---event-date specification and spread definition---producing six event-study configurations in total: three Korea event-date variants (narrow 2024 rollout, window sensitivity, and spaced follow-through) each estimated against both the primary MSCI EM Asia and the secondary TOPIX spreads.

\paragraph{Window sensitivity.} The window sensitivity variant truncates the post-event horizon to 12 months, testing whether the primary results depend on the longer 20-month window. The 12-month CARs are numerically identical to those in the primary specification by construction. For the primary MSCI EM Asia spread, the 12-month CARs are $-2.178$ (launch), $-1.804$ (guidelines), and $-2.121$ (implementation push)---all negative and consistent with the primary result. For the secondary TOPIX spread, the 12-month CARs are $-0.289$ (launch), $+1.030$ (guidelines), and $+1.426$ (implementation push), replicating the confound pattern discussed above. The primary MSCI EM Asia finding is therefore robust to the truncated window. Figure~6 presents both panels of the window-sensitivity specification.

\begin{figure}[htbp]
    \centering
    \begin{subfigure}[b]{0.49\textwidth}
        \centering
        \includegraphics[width=\textwidth]{figure_korea_asia_event_study_robustness_narrow_2024_rollout_post12.png}
        \caption{KOSPI--MSCI EM Asia spread, 12-month window (primary)}
        \label{fig:rob_post12_asia}
    \end{subfigure}
    \hfill
    \begin{subfigure}[b]{0.49\textwidth}
        \centering
        \includegraphics[width=\textwidth]{figure_korea_event_study_robustness_narrow_2024_rollout_post12.pdf}
        \caption{KOSPI--TOPIX spread, 12-month window (secondary)}
        \label{fig:rob_post12_topix}
    \end{subfigure}
    \caption{Window-sensitivity robustness: Korea Value-Up event study truncated to a 12-month post-event window. CARs at $\tau = 12$ are identical to the primary specification (Figures~4 and 5) by construction. Panel~(a) shows the primary MSCI EM Asia spread; Panel~(b) shows the secondary TOPIX spread. All three primary MSCI EM Asia CARs remain negative at the 12-month horizon, confirming that the result does not depend on the extended 20-month window.}
    \label{fig:rob_post12}
\end{figure}

\paragraph{Spaced follow-through.} Figure~7 presents the spaced follow-through specification, which replaces the clustered 2024 dates with a temporally spread event sequence: the February 2024 program launch, the July 2025 mandatory governance disclosure expansion, and the February 2026 Value-Up dividend disclosure rule. This reduces calendar-time overlap at the cost of a 2-month post-event window for the most recent events given the April 2026 study endpoint. For the program launch, the primary MSCI EM Asia 2-month CAR is $-0.468$ and the secondary TOPIX CAR is $-0.362$---both negative and consistent with the primary result. For the July 2025 mandatory expansion, the MSCI EM Asia CAR is $-0.006$ (near-zero) and the TOPIX CAR is $+0.654$; for the February 2026 dividend disclosure rule, the MSCI EM Asia CAR is $+1.324$ and the TOPIX CAR is $+1.825$. The two most recent mandatory events show tentatively positive short-run CARs in both specifications, but a 2-month window is insufficient to distinguish a sustained structural improvement from short-run announcement effects. We report these results for transparency; the spaced follow-through design was constructed to reduce overlap rather than to generate inference about the very latest legislative milestones.

\begin{figure}[htbp]
    \centering
    \begin{subfigure}[b]{0.49\textwidth}
        \centering
        \includegraphics[width=\textwidth]{figure_korea_asia_event_study_robustness_spaced_follow_through.png}
        \caption{KOSPI--MSCI EM Asia spread (primary)}
        \label{fig:rob_spaced_asia}
    \end{subfigure}
    \hfill
    \begin{subfigure}[b]{0.49\textwidth}
        \centering
        \includegraphics[width=\textwidth]{figure_korea_event_study_robustness_spaced_follow_through.pdf}
        \caption{KOSPI--TOPIX spread (secondary)}
        \label{fig:rob_spaced_topix}
    \end{subfigure}
    \caption{Spaced follow-through robustness: event sequence replaced with the February 2024 program launch, the July 2025 mandatory governance disclosure expansion, and the February 2026 Value-Up dividend disclosure rule. Post-event windows are 2 months for the two most recent events given the April 2026 study endpoint. Panel~(a) uses the primary MSCI EM Asia spread; Panel~(b) uses the secondary TOPIX spread. The launch cohort replicates the negative primary result ($-0.468$ MSCI EM Asia, $-0.362$ TOPIX at $\tau = 2$). The very short window for the 2025 and 2026 events precludes durable inference.}
    \label{fig:rob_spaced}
\end{figure}

Across all six robustness configurations, the primary MSCI EM Asia conclusion is stable. The three primary 2024 Value-Up milestones produce uniformly negative 12-month CARs of $-2.18$, $-1.80$, and $-2.12$ P/B points in the MSCI EM Asia specification, unaffected by TOPIX-specific dynamics. The secondary TOPIX specification shows positive CARs for the two later milestones, consistent with the TOPIX-confound explanation and not with genuine Korean valuation improvement. These results, combined with the Japan benchmark evidence that Japan's 2023 hard mandate produced a $-2.58$ P/B point CAR in the primary MSCI EM Asia--TOPIX specification (and $-6.48$ in the secondary KOSPI--TOPIX specification), jointly support the central claim: Korea's voluntary Corporate Value-Up program has not generated a structural improvement in Korea's relative valuation, and the governance-reform gap widened precisely when Japan escalated from soft codes to hard exchange-level mandates while Korea did not.

\section{Discussion}

\subsection{Interpreting the Korea--Japan Divergence}

The central empirical finding is a sharp asymmetry in the valuation consequences of Japan's three governance reform milestones. In the primary MSCI EM Asia--TOPIX specification, Japan's two early soft-law interventions---the 2014 Stewardship Code and the 2015 Corporate Governance Code---produced positive abnormal spread movements, indicating that the entire EM Asia regional basket held its position against Japan in the immediate post-reform periods. The 2023 TSE P/B Reform Request is categorically different: the MSCI EM Asia--TOPIX spread declined monotonically to $-2.584$ P/B points by $\tau = 20$, and the secondary KOSPI--TOPIX spread widened to $-6.481$ P/B points by $\tau = 24$---the largest event in the sample. Japan's mandatory P/B reform lifted TOPIX above the entire regional peer basket, with Korea's individual shortfall exceeding the regional average.

This asymmetry carries a specific interpretation. Soft governance codes---whether Japan's 2014 Stewardship Code or its 2015 Corporate Governance Code---are instruments of persuasion: they articulate norms for institutional investor engagement and board structure but impose no sanctions for non-compliance. Under voluntary frameworks, firms with the greatest incentives to resist reform---those in which controlling shareholders extract private benefits through related-party transactions---are precisely the ones least likely to opt in. The empirical record from Korea's own institutional-activism literature corroborates this: markets do not react to NPS vote-no announcements in the short run \citep{kim_nps_governance_firm_value}, suggesting that voluntary shareholder activism lacks the enforcement credibility necessary to shift valuation expectations. Japan's two early soft-code events therefore did not generate a sustained TOPIX re-rating that Korea failed to match, because neither code constituted an enforceable governance shock for the firms most responsible for Japan's pre-reform discount.

The 2023 TSE mandate crossed a qualitative threshold. By requiring firms with P/B below 1.0 to publicly disclose capital efficiency improvement plans or explain non-compliance---backed by the threat of delisting guidance and active TSE follow-up---the mandate transformed governance reform from a reputational norm into a compliance obligation with observable consequences. The immediate negative response in the MSCI EM Asia--TOPIX spread reflects investors anticipating differential capital allocation improvement: Japan's listed universe would improve while the broader EM Asia region, Korea included, would not. Korea's 2024 Value-Up program generated no comparable improvement: against the primary MSCI EM Asia benchmark, all three milestones produced persistently negative CARs of $-2.18$, $-1.80$, and $-2.12$ P/B points at the 12-month horizon---Korea was losing ground to regional peers who had themselves failed to keep pace with Japan.

The contrast between Japan's hard mandate and Korea's voluntary program is therefore not merely a question of intensity but of mechanism. Voluntary programs do not resolve the fundamental collective-action problem: controlling shareholders who benefit from the status quo have no individually rational reason to comply, and markets will not price in a reform they do not expect to be enforced. Whether the divergence---$-2.58$ P/B points in the primary MSCI EM Asia specification, $-6.48$ in the secondary KOSPI--TOPIX specification---fully closes the structural gap in the long run will depend on Japan maintaining enforcement and on Korea's eventual transition to comparable mandates.

\subsection{The Governance Channel and Its Limits}

The results of this paper are consistent with, but do not definitively establish, a governance channel as the exclusive mechanism behind the Korea Discount. Three alternative interpretations merit explicit consideration.

First, the discount could partly reflect structural differences in the industrial composition of KOSPI relative to TOPIX and MSCI EM Asia. Korean equities are heavily weighted toward manufacturing conglomerates with high tangible-asset intensity, precisely the firm characteristics identified as correlates of low P/B in cross-sectional regressions \citep{yang_korea_discount_growth_vs_governance}. Our index-level analysis cannot control for compositional differences of this kind, and the sub-period deterioration in Korea's mean P/B from 1.36 (pre-reform) to 0.93 (post-TSE) does not by itself isolate governance from slower structural transitions in Korean industry. However, the event-study design partially addresses this concern: if compositional differences are approximately stable within the 36-month pre-event window used for trend estimation, then the abnormal spread in the post-event period captures reform-related deviations rather than secular composition shifts. The consistency of negative CARs across three distinct Korea Value-Up milestones, spanning February through August 2024, is harder to explain through slow-moving compositional change than through a market-level assessment of Korea's governance trajectory.

Second, foreign exchange dynamics could conflate valuation and currency effects if P/B ratios are being compared across currency boundaries in a way that MSCI index methodology does not fully equalize. The panel OLS results partially address this concern: the log FX coefficient in the two-way fixed effects specification is $-0.024$ ($p = 0.018$), indicating a modest negative association between won depreciation and Korean P/B after controlling for country and time fixed effects, but the reform interaction terms for the TSE Corporate Governance Code and the TSE P/B Reform Request retain significance ($-0.322$, $p < 0.001$ and $-0.234$, $p = 0.056$ respectively), and the event-study design absorbs FX trends through the pre-event linear extrapolation. Currency effects alone cannot account for the persistence and magnitude of the divergence documented here.

Third, the growth-versus-governance debate in the literature---cross-sectional P/B regressions find that R\&D investment and firm maturity outperform governance variables \citep{yang_korea_discount_growth_vs_governance, lee_korea_discount_pbr}---raises the question of whether governance reform can raise valuations without also addressing the underlying growth deficits of large Korean conglomerates. Our results do not directly engage this question at the firm level, since we use index-level data. What the event-study evidence does establish is that Japan's market reacted positively to a governance mandate even though Japanese conglomerates share many of the same structural growth limitations as their Korean counterparts. If governance improvements were insufficient to raise valuations in the absence of fundamental growth, the TSE mandate should not have produced a sustained TOPIX re-rating. The magnitude of the 2023 CAR is therefore at least partially inconsistent with a pure growth story, and more consistent with the view that governance-induced capital allocation improvements---including buybacks, dividend increases, and divestiture of low-return subsidiaries---can drive re-rating independently of underlying business model transformation.

\subsection{Limitations}

Several limitations of the empirical design should be acknowledged. First, the stacked event-study methodology produces descriptive CARs rather than inferential estimates: the saturated cohort-by-relative-time specification places one observation per cell, precluding the construction of standard errors, and the clustering of Korea's three Value-Up milestones within a seven-month window creates calendar-time overlap that violates the standard stacked-DiD independence assumption. The CARs reported in Section~4 are therefore properly interpreted as measures of realized abnormal spread movement, not as causal effect estimates with associated uncertainty. Readers should be cautious about reading precision into the specific numerical values; the qualitative patterns---persistent negation in the MSCI EM Asia comparison, the sharp sign reversal between Japan's soft and hard codes---are robust to the design's descriptive limitations.

Second, Japan serves as a governance reform benchmark rather than a clean causal counterfactual. The identifying assumption underlying the primary Japan event study is that, absent Japan's governance reforms, the MSCI EM Asia--TOPIX spread would have continued along its pre-event linear trend. This assumption is implausible for long horizons---the indices respond to different macroeconomic shocks, exchange rate regimes, and sector compositions---but defensible within the 36-month estimation window used here. The secondary KOSPI--TOPIX analysis carries the same assumption for Korea specifically; the convergence of the primary EM Asia and secondary KOSPI results across all three reform events strengthens the inference that Japan's hard mandate, not Korea-idiosyncratic dynamics, drove the divergence.

Third, the analysis is conducted entirely at the index level and cannot speak to the cross-sectional distribution of governance improvements within Korea. The market response to Korea's 2025 Commercial Act amendment was heterogeneous: firms with the highest book-to-market ratios and the lowest controlling-shareholder ownership concentrations showed the largest abnormal returns, consistent with governance reform having its greatest effect precisely where expropriation risk is highest \citep{kang_commercial_act_market_reactions}. Our aggregate index-level design would not detect such heterogeneity even if it were present in the data, and the apparent absence of aggregate improvement in the KOSPI--MSCI EM Asia spread does not preclude meaningful governance gains at the firm level for a subset of Korean companies. Firm-level event studies around Korea's 2024 and 2025 reform milestones, using within-firm variation in governance exposure as a conditioning variable, would be a natural complement to the aggregate evidence presented here.

Fourth, the study window ends in April 2026, and only two post-event months are available for Korea's most recent mandatory reform events. The spaced follow-through specification suggests tentatively positive CARs for the July 2025 mandatory governance disclosure expansion and the February 2026 Value-Up dividend disclosure rule, but a 2-month window is entirely insufficient to distinguish a sustained structural improvement from short-run announcement effects. The conclusions of this paper pertain to the voluntary 2024 Value-Up phase; whether Korea's subsequent mandatory increments---and the broader 2025 Commercial Act fiduciary duty expansion---produce durable re-rating is an empirical question that cannot yet be answered with the available data.

\subsection{Policy Implications}

The evidence in this paper, read alongside the literature surveyed in Section~2, supports a specific policy conclusion: resolving the Korea Discount requires hard legislative reform that removes the voluntary character of governance compliance, rather than iterative enhancement of existing soft-law frameworks. The mechanism through which hard mandates operate is not primarily informational---markets are not learning new facts about Korean corporate governance from a mandatory disclosure rule that they could not have inferred from voluntary disclosures. Rather, mandatory frameworks change the strategic environment for controlling shareholders by making non-compliance costly and observable, thereby credibly committing the regime to enforcement in a way that voluntary guidelines cannot.

This distinction has direct relevance for Korea's reform agenda. The 2024 Corporate Value-Up program failed not because its objectives were poorly specified but because the voluntary mechanism left the participation decision with the very firms whose controlling shareholders had the largest private incentives to resist. A hard mandate analogous to the TSE P/B directive---requiring firms below book value to disclose capital efficiency improvement plans under threat of regulatory consequence---would remove that opt-out. The 2025 Korean Commercial Act amendment, which expanded directors' fiduciary duties from the corporation to shareholders broadly, represents meaningful progress: legislative passage produced significant positive abnormal returns for high-book-to-market firms, consistent with markets revising expropriation risk downward \citep{kang_commercial_act_market_reactions}. However, a fiduciary duty expansion creates a legal cause of action rather than a proactive compliance obligation; whether it shifts controlling shareholder behavior without accompanying exchange-level enforcement remains to be seen.

Governance reform and market openness are complementary instruments \citep{joo_foreign_investors_value_up}: foreign institutional investors impose discipline on both over-investing and under-returning firms, correcting the allocation distortions that flow from concentrated ownership structures. Reforms that reduce the Korea Discount through improved governance would therefore be self-reinforcing if they simultaneously attract foreign institutional participation, because greater foreign ownership raises the share of investors who monitor management aggressively and vote against related-party transactions. The Korean government's parallel efforts to improve market accessibility---including English-language corporate disclosures and simplified foreign investor registration---should be understood in this context as part of a governance reform package rather than a standalone capital-account liberalization measure.

The findings also bear on the broader debate over the relative contribution of governance failures and growth deficits to the Korea Discount. The evidence from Japan's 2023 mandate suggests that governance reforms can produce large, sustained valuation improvements even within an industrial structure broadly similar to Korea's---without requiring a transformation of the underlying business models of the large conglomerates that dominate both markets. This does not imply that growth-oriented reforms are irrelevant, but it does establish a lower bound on the valuation impact that governance reform alone can achieve. Policies that simultaneously address governance---through mandatory disclosure obligations, enhanced fiduciary duties, and exchange-level enforcement---and growth---through R\&D tax incentives, cross-shareholding reduction, and holding company discount compression---are likely to be complementary and mutually reinforcing, as the literature on Korea's chaebol structure suggests that governance improvements reduce value-destroying pyramidal acquisitions that otherwise crowd out productive investment.

\subsection{Conclusion}

This paper has argued and empirically documented that the Korea Discount is structural, governance-driven, and unlikely to close without hard legislative reform. The descriptive panel establishes KOSPI's mean discount of $-0.177$ P/B points to TOPIX and $-0.601$ to MSCI EM over 20 years, with a sub-period deterioration from near-parity in 2004--2013 to a 0.43 P/B point gap in 2023--2024. The geopolitical risk falsification rules out the North Korean threat explanation ($t = -0.835$, $p = 0.405$). The Japan event study identifies the 2023 TSE hard mandate as the structural inflection point: the primary MSCI EM Asia--TOPIX CAR reached $-2.58$ P/B points at $\tau = 20$, and the secondary KOSPI--TOPIX CAR reached $-6.48$ at $\tau = 24$---the largest event in the sample---while Japan's earlier soft codes produced positive CARs in both specifications. The Korea Value-Up event study shows the voluntary program produced no analogous improvement: the primary MSCI EM Asia specification records persistently negative CARs of $-2.18$, $-1.80$, and $-2.12$ P/B points at $\tau = 12$ for all three milestones, robust across event-date specifications and post-event window lengths.

The practical implication is that Korea currently stands where Japan stood before 2023: in possession of well-articulated soft-law governance norms whose voluntary character renders them structurally insufficient to overcome the incentive problems of concentrated ownership. Japan's experience demonstrates that the transition to hard exchange-level mandates can produce large, rapid, and sustained valuation improvements---improvements that were visible not just in Korea's relative performance but across the entire EM Asia regional basket. The 2025 Commercial Act amendment expanded the legal framework for shareholder protection in a way that markets have already partially priced, but an exchange-level enforcement mechanism---requiring firms below book value to disclose credible capital efficiency plans under threat of regulatory consequence---remains the missing complement. Until that transition occurs, Korea's valuation gap will persist as a structural equilibrium outcome of an ownership regime that voluntary guidelines have repeatedly failed to reform.


\bibliographystyle{plainnat}
\bibliography{citations}

\end{document}
